\section{Data and baselines}
\label{sec:baselines}

We aim to change one CMB lens while holding everything else as fixed as possible. Throughout, $A$ denotes the SPT-3G 2018 lens and $B$ the SPT-3G D1 lens. We write $\dchi_{\mathrm{D1}|2018}$ for $\dchi_{B|A}$ (D1 tested at the 2018 fit) and $\dchi_{2018|\mathrm{D1}}$ for the reverse.

\subsection{The two SPT-3G lenses}
\label{sec:spt-lenses}

Both lenses are multifrequency TT/TE/EE likelihoods evaluated with \texttt{candl} \cite{candl}:
\begin{itemize}[leftmargin=*, itemsep=2pt]
  \item \textbf{SPT 2018} \cite{spt3g2018_ttteee} (\texttt{candl\_data.SPT3G\_2018\_TTTEEE\_multifreq}): TT for $750\le\ell\le3000$, TE and EE for $300\le\ell\le3000$.
  \item \textbf{SPT D1} \cite{spt3g_d1} (\texttt{spt\_candl\_data.SPT3G\_D1\_TnE\_multifreq}): TT for $400\le\ell\le3000$, TE and EE for $400\le\ell\le4000$.
\end{itemize}
These ranges follow \cite{arxiv2601_16277}. Each release carries its own nuisance model and internal priors, and we keep them as part of the lens. In particular, the two releases include different Gaussian priors on the optical depth $\tau$: mean $0.054$ and standard deviation $0.0074$ for 2018, and mean $0.051$ and standard deviation $0.006$ for D1. A lens swap therefore changes this prior as well as the data. \Cref{sec:results-audits} tests how much this matters.

\subsection{Planck subset and lensing}
\label{sec:planck-subset}

To mirror the CMB baseline of \cite{arxiv2601_16277}, we add a subset of Planck 2018 (PR3) products and Planck PR4 lensing \cite{planck2018_likelihood,planck_pr4_lensing,planck_pr4_lensing_maps}:
\begin{itemize}[leftmargin=*, itemsep=2pt]
  \item low-$\ell$ TT and EE for $\ell<30$ (EE via the \texttt{EE\_sroll2} variant in Cobaya);
  \item high-$\ell$ TT restricted to $30\le\ell\le756$ (\texttt{planck\_2018\_highl\_plik.\allowbreak TT\_lite\_native} with \texttt{bins\_for\_L\_range=\allowbreak[30,756]});
  \item no Planck high-$\ell$ TE or EE;
  \item Planck PR4 lensing (\texttt{planckpr4lensing.PlanckPR4Lensing}).
\end{itemize}
This is not the full Planck 2018 likelihood \cite{planck2018_cosmoparams}. Unlike \cite{arxiv2601_16277}, we include no SPT lensing likelihood: neither the 2018 lensing spectrum nor the D1-era MUSE lensing spectrum that they swap in alongside D1. Our ``full CMB'' baseline therefore swaps only the SPT TT/TE/EE likelihood. \Cref{tab:baseline-defs} compares the two definitions.

\begin{table}[t]
  \centering
  \caption{CMB baselines as defined in \cite{arxiv2601_16277}. Our ``full CMB'' baseline uses the same SPT primary multipole ranges and Planck subset, keeps Planck PR4 lensing, and omits both SPT lensing likelihoods, so that only the SPT TT/TE/EE lens changes.}
  \label{tab:baseline-defs}
  \small
  \begin{tabular}{@{} p{0.14\linewidth} p{0.22\linewidth} p{0.31\linewidth} p{0.25\linewidth} @{}}
\toprule
Baseline & SPT primary ($\ell$ ranges) & Planck PR3 primary subset & Lensing \\
\midrule
CMB & TT $[750,3000]$; TE/EE $[300,3000]$ & TT $\ell<30$ (Commander); EE $\ell<30$ (SRoll2); TT $30<\ell<756$ (Plik) & SPT-3G 2018 lensing + Planck PR4 lensing \\
\addlinespace
CMB-D1 & TT $[400,3000]$; TE/EE $[400,4000]$ & as above & SPT-3G 2019--2020 MUSE lensing + Planck PR4 lensing \\
\bottomrule
\end{tabular}

\end{table}

\subsection{DESI DR2 BAO}
\label{sec:desi-bao}

For BAO we use a Gaussian likelihood built from the public DESI DR2 BAO mean vector and covariance for the combined ``ALL'' sample \cite{desi_dr2_bao_docs,bao_data_v26,desi_dr2_results_ii_bao,desi_dr2_results_i_lya}. It includes $D_M/r_d$, $D_H/r_d$ and $D_V/r_d$ points from galaxy and Lyman-$\alpha$ forest BAO \cite{eisenstein2005_baopeak}. Here $D_M=(1+z)D_A$, $D_H=c/H(z)$ and $D_V=(zD_HD_M^2)^{1/3}$. Predictions use the same Boltzmann code as the CMB likelihoods in each run.

\subsection{Two audit baselines}
The audits and posteriors use two completions:
\begin{itemize}[leftmargin=*, itemsep=2pt]
  \item \textbf{SPT+DESI}: one SPT lens plus DESI DR2 BAO, computed with CAMB.
  \item \textbf{Full CMB+DESI}: one SPT lens plus the Planck subset, Planck PR4 lensing and DESI DR2 BAO, computed with CLASS.
\end{itemize}
Both use $\Lambda$CDM$+\mnu$ with three degenerate massive neutrinos, a flat prior on $\mnu\in[0,5]$\,eV, and the six standard $\Lambda$CDM parameters. The two baselines also differ in Boltzmann code and in which foreground parameters are free. A difference between them therefore cannot be attributed to the Planck subset alone.
